\documentclass[11pt]{amsart}
\usepackage{amsmath,amssymb,amsthm}
\usepackage[margin=1.05in]{geometry}

\numberwithin{equation}{section}
\newtheorem{theorem}{Theorem}[section]
\newtheorem{proposition}[theorem]{Proposition}
\newtheorem{lemma}[theorem]{Lemma}
\newtheorem{corollary}[theorem]{Corollary}
\theoremstyle{definition}
\newtheorem{definition}[theorem]{Definition}
\theoremstyle{remark}
\newtheorem{remark}[theorem]{Remark}
\newtheorem{example}[theorem]{Example}

\newcommand{\Z}{\mathbb{Z}}
\newcommand{\Q}{\mathbb{Q}}
\newcommand{\R}{\mathbb{R}}
\newcommand{\N}{\mathbb{N}}
\newcommand{\C}{\mathbb{C}}
\newcommand{\F}{\mathbb{F}}
\newcommand{\cA}{\mathcal{A}}
\newcommand{\cP}{\mathcal{P}}
\newcommand{\pp}{\mathfrak{p}}
\newcommand{\ff}{\mathfrak{f}}
\newcommand{\Ns}{\mathrm{N}}
\newcommand{\Cl}{\operatorname{Cl}}
\newcommand{\Div}{\operatorname{Div}}
\newcommand{\Frob}{\operatorname{Frob}}
\newcommand{\Gal}{\operatorname{Gal}}
\newcommand{\Prob}{\operatorname{Prob}}
\newcommand{\Sp}{\operatorname{Sp}}
\newcommand{\id}{\mathrm{id}}

\PassOptionsToPackage{hyphens}{url}
\usepackage[hidelinks,bookmarks=false]{hyperref}

\title[Bost--Connes systems over function fields]{Localized Bost--Connes systems over global function fields:\\ commensurability forbids type $\mathrm{III}_1$,\\ and the factor recovers the temperature}
\author{Ruqing Chen}
\address{GUT Geoservice Inc., Montreal, Canada}
\email{ruqing@hotmail.com}
\thanks{Sources, code and data for the entire series: \url{https://github.com/Ruqing1963/localized-bost-connes}; this paper is in \texttt{papers/XVII-function-fields/}.}
\date{}

\begin{document}

\begin{abstract}
We transport the localized Bost--Connes theory to $K=\F_q(C)$ for a smooth projective
geometrically irreducible curve $C/\F_q$ of genus $g$, and find that one structural feature
inverts and, with it, the central negative result of Paper IX of the series.

Over a number field the logarithms $\log\Ns\pp$ of an infinite prime set are never pairwise
commensurable \cite[Lem.~4.1]{SeriesI}; over a function field they are \emph{always}
commensurable, being $\deg(P)\log q$. Hence
$\Gamma_\beta=\langle q^{-\beta\deg P}\rangle=q^{-\beta d_S\Z}$ with
$d_S=\gcd_{P\in S}\deg P$ is a \emph{discrete} subgroup of $\R_{>0}$, so
$S(M)\subseteq\Gamma_\beta\cup\{0\}$ can never be $[0,\infty)$: \textbf{the extremal
$\mathrm{KMS}_\beta$ factors are never of type $\mathrm{III}_1$}, in complete contrast with
\cite[Thm.~3.2]{SeriesI}, where $\mathrm{III}_1$ is generic. The proposed alternative
``$\mathrm{III}_0$ if $\Gamma_\beta$ is dense'' cannot occur, a discrete group not being
dense; the alternatives are $\mathrm{III}_\lambda$ with $\lambda\in q^{-\beta d_S\Z}$,
$\mathrm{III}_0$, or semifiniteness.

The consequence is the interesting one. In the $\mathrm{III}_\lambda$ case
$M_\beta\cong R_\lambda$ with $\lambda=q^{-\beta d_S}$, and the Powers factors are mutually
non-isomorphic; so \textbf{the von Neumann algebra determines $\beta\,d_S\log q$}. Paper IX
showed that over number fields every system with $S=P_K$, $\beta\le1$ gives $R_\infty$ and the
factor remembers neither $K$, nor $S$, nor $\beta$. Over function fields the collapse is
partial: the temperature survives weak closure. This is the first repair of Theorem 2.1 of
Paper IX in the whole programme, and it is bought by commensurability rather
than by any rigidity input.

Two corrections. The modular flow is periodic with fundamental period
$T_0=2\pi/(d_S\log q)$, not $2\pi/\log q$; the two agree exactly when $d_S=1$, which does hold
for $S=\cP_K$ by F.~K. Schmidt's theorem but fails for, say, the set of points of even degree.
And in the Weil--Chebotarev bound the term $(g_L+1)q^{n/2}$ swamps the main term
$2g_Lq^{n/2}/n$ and is superfluous: the truth is $O_\chi(q^{n/2}/n)$, with the
constant-field-extension characters requiring a separate and elementary treatment, Weil's
bound not applying to them.

The rest goes through and is unconditional. $\beta_c=1$ for every curve; uniqueness holds on
$(0,1]$ for $S=\cP_K$, unconditionally --- the class-group characters, which over a number
field can obstruct permanently, here have detecting sets of positive density; and a set with
$A_n(S)\asymp q^n/n^2$ has $\sum_{P\in S}q^{-\deg P}=\zeta(2)$-like and finite, giving a
\emph{closed} low-temperature phase $[1,\infty)$ with no sieve hypothesis whatever.
\end{abstract}

\maketitle

\section{The system and its periodic flow}

Let $C/\F_q$ be a smooth projective geometrically irreducible curve of genus $g$,
$K=\F_q(C)$, and $\cP_K$ the set of closed points. For $S\subseteq\cP_K$ put
\[
J_S=\Div_+^S(C)=\bigoplus_{P\in S}\N\cdot P,\qquad
\Ns(P)=q^{\deg P},\qquad
\cA_{C,S}=C(Y_S)\rtimes J_S,
\]
with $G_S=\mathbb I_S/\overline{K^\times_S}$ and $\sigma_t(\mu_D)=q^{it\deg D}\mu_D$. Write
\[
d_S:=\gcd_{P\in S}\deg P .
\]
Constructions and notation follow \cite{Main}; Papers VIII, IX, XIII and XV of the series,
referred to by numeral, are unpublished companions, and every fact used from them is
restated or verified in place. One convention is needed at once: for $S=\cP_K$ the degree
map gives $\mathbb I_S/\overline{K^\times_S}$ a quotient $\Z$, so the group as written is
not compact; as the framework of \cite{Main} requires a compact symmetry group (and as is
standard for Bost--Connes systems over function fields), $G_S$ denotes throughout the
profinite completion in the degree direction. Its dual $\widehat{G_S}$ is then a torsion
group, and the characters factoring through the constant-field tower
$\Gal(\overline\F_q/\F_q)\cong\widehat\Z$ are exactly $P\mapsto\zeta^{\deg P}$ with $\zeta$
a root of unity. For the basic theory of curves over finite fields we refer to
\cite{Rosen,Stichtenoth}.

\begin{proposition}[Periodicity]\label{prop:period}
The flow $\sigma$ is periodic with fundamental period
\[
T_0=\frac{2\pi}{d_S\log q}.
\]
In particular $T_0=2\pi/\log q$ precisely when $d_S=1$, which holds for $S=\cP_K$ by
F.~K. Schmidt's theorem \cite{Rosen} that the gcd of the degrees of closed points of a curve over a finite
field is $1$; it fails for instance when $S$ consists of the points of even degree, where
$d_S=2$ and the period is halved.
\end{proposition}

\begin{proof}
$\sigma_T=\id$ iff $q^{iT\deg D}=1$ for every $D\in J_S$, i.e.\ $T\log q\deg D\in2\pi\Z$. The
degrees realized by $J_S$ generate the subgroup $d_S\Z$ of $\Z$, so the condition is
$T\,d_S\log q\in2\pi\Z$, whose least positive solution is $T_0$.
\end{proof}

\begin{remark}\label{rem:periodic}
Periodicity of $\sigma$ has no analogue over number fields, where the $\log\Ns\pp$ generate a
dense subgroup of $\R$ and $\sigma$ is never periodic. It means that $\sigma$ factors through
the circle $\R/T_0\Z$, so $\cA_{C,S}$ carries a genuine $\mathbb T$-action rather than a flow.
\end{remark}

\section{Rational thermodynamics}

\begin{proposition}[Weil zeta]\label{prop:zeta}
With $T=q^{-\beta}$,
\[
\prod_{P\in S}\bigl(1-q^{-\beta\deg P}\bigr)^{-1}=Z_S(C,T),
\]
and for $S=\cP_K$ this is the Weil zeta function \cite{Weil}
$Z_C(T)=\dfrac{P(T)}{(1-T)(1-qT)}$ with $P(T)=\prod_{j=1}^{2g}(1-\alpha_jT)$ and
$|\alpha_j|=q^{1/2}$. Its poles are at $T=1$ and $T=q^{-1}$, that is at $\beta=0$ and
$\beta=1$.
\end{proposition}

\begin{corollary}[Critical exponent]\label{cor:betac}
For every curve $C$ and $S=\cP_K$, $\beta_c(C)=1$: the series
$\sum_{P}q^{-\beta\deg P}=\sum_n a_nq^{-\beta n}$, with
$a_n=\#\{P:\deg P=n\}=\frac{q^n}{n}+O\bigl(\frac{q^{n/2}}{n}\bigr)$ by Weil, converges for
$\beta>1$ and diverges for $\beta\le1$.
\end{corollary}

\begin{proof}
$\sum_n\frac{q^{n(1-\beta)}}{n}$ diverges for $\beta\le1$, converging otherwise; the Weil
error term contributes $\sum_nq^{n(1/2-\beta)}/n$, convergent for $\beta>1/2$.
\end{proof}

\begin{remark}[Numerical illustration]\label{rem:numbetac}
Partial sums $\sum_{n\le40}a_nq^{-\beta n}$ with $a_n=q^n/n$: for $q=2$ they are $11.36$,
$4.28$, $2.69$ at $\beta=0.9,1.0,1.1$, and for $q=4$ they are $65.68$, $4.28$, $2.04$; the
value at $\beta=1$ is $\sum_{n\le40}1/n$ in both cases, growing logarithmically, and the
$\beta=1.1$ column is bounded.
\end{remark}

\section{Commensurability forbids type $\mathrm{III}_1$}

\begin{lemma}[Total commensurability]\label{lem:commens}
For every $P\in S$, $\log\Ns(P)=\deg(P)\log q\in(\log q)\Z$. Hence the log-norms are
pairwise commensurable, and
\[
\Gamma_\beta=\bigl\langle\Ns(P)^{-\beta}:P\in S\bigr\rangle=q^{-\beta d_S\Z},
\]
a discrete subgroup of $(\R_{>0},\times)$.
\end{lemma}

\begin{remark}[The inversion]\label{rem:inversion}
\cite[Lem.~4.1]{SeriesI} proves that over a number field an infinite $S$ never has pairwise
commensurable log-norms, distinct rational primes having $\Q$-independent logarithms. Over a
function field the opposite holds universally. This single reversal drives everything below.
\end{remark}

\begin{theorem}[Never $\mathrm{III}_1$]\label{thm:neverIII1}
For every curve $C$, every $S$ and every $\beta>0$, the extremal $\mathrm{KMS}_\beta$ factors
$M=\pi_\varphi(\cA_{C,S})''$ are \emph{not} of type $\mathrm{III}_1$. Moreover
$\Gamma_\beta$ being discrete, the case ``$\Gamma_\beta$ dense'' cannot occur, so the
possibilities are: type $\mathrm{III}_\lambda$ with $\lambda\in q^{-\beta d_S\Z}$; type
$\mathrm{III}_0$, where $S(M)=\{0,1\}$; or semifinite, where $S(M)=\{1\}$.
\end{theorem}

\begin{proof}
By \cite[Prop.~6.1]{Main}, $S(M)\subseteq\overline{\Gamma_\beta}\cup\{0\}$, and
$\overline{\Gamma_\beta}=\Gamma_\beta$ by Lemma~\ref{lem:commens}. Type $\mathrm{III}_1$
requires $S(M)=[0,\infty)$, which is not contained in a discrete set together with $0$. The
listed cases exhaust Connes' classification of $S(M)$.
\end{proof}

\begin{corollary}[Contrast]\label{cor:contrast}
By \cite[Thm.~3.2]{SeriesI}, over a number field with $S=P_K$ and $0<\beta\le1$ the type is
always $\mathrm{III}_1$. Over a function field it never is, for any $S$ and any $\beta$. The
two settings are opposite in this respect, and the reason is Lemma~\ref{lem:commens} alone,
no arithmetic being used.
\end{corollary}

\section{The factor recovers the temperature}

\begin{theorem}[Partial repair of Paper IX]\label{thm:repair}
Suppose $M_\beta=\pi_{\varphi_\beta}(\cA_{C,S})''$ is of type $\mathrm{III}_\lambda$ with
$\lambda=q^{-\beta d_S}$. Then $M_\beta\cong R_\lambda$, the Powers factor: $\cA_{C,S}$ is
amenable --- the argument of Paper IX (its Lemma 1.1) applies verbatim, the acting group
$I_S=\Z^{(S)}$ being abelian here too, so the groupoid is amenable and $M_\beta$ is
injective --- and since $R_\lambda\cong R_{\lambda'}$ only for
$\lambda=\lambda'$, the isomorphism class of $M_\beta$ determines
\[
\lambda=q^{-\beta d_S},\qquad\text{hence}\qquad \beta\,d_S\log q .
\]
In particular, for fixed $q$ and $d_S$, distinct temperatures give non-isomorphic factors.
\end{theorem}

\begin{proof}
Amenability and Connes' classification give $M_\beta\cong R_\lambda$; the Powers factors are
pairwise non-isomorphic, their $S$-invariants $\{0\}\cup\lambda^\Z$ being distinct.
\end{proof}

\begin{remark}[What this changes]\label{rem:changes}
Theorem 2.1 of Paper IX showed that over number fields
$M_{\Q,P_\Q,1}\cong M_{\Q,P_\Q,1/2}\cong M_{K,P_K,\beta}\cong R_\infty$ for every $K$ and
every $\beta\in(0,1]$: the factor remembers neither the field, nor the prime set, nor the
temperature. Theorem~\ref{thm:repair} is the first place in this programme where that
collapse is not total. It is worth being precise about how much is repaired: the factor
recovers the single real number $\beta d_S\log q$, and nothing else --- not $C$, not $g$, not
$S$ beyond $d_S$, not the obstruction group. But $\beta$ was the quantity whose loss was most
striking in Paper IX, and it is recovered here for free, by commensurability rather
than by any rigidity theorem.

Numerically:
\[
\begin{array}{lccccc}
(q,d_S,\beta) & (2,1,\tfrac12) & (2,1,1) & (2,1,2) & (4,2,1) & (9,1,\tfrac32)\\
\lambda & 0.7071 & 0.5000 & 0.2500 & 0.0625 & 0.0370
\end{array}
\]
Distinct temperatures give distinct factors.
\end{remark}

\begin{remark}[The remaining ambiguity]\label{rem:ambiguity}
Only the product $\beta d_S\log q$ is recovered, so $(q,d_S,\beta)$ and
$(q',d_S',\beta')$ with $\beta d_S\log q=\beta'd_S'\log q'$ give isomorphic factors: for
example $(q,d_S,\beta)=(4,1,1)$ and $(2,1,2)$ both give $\lambda=1/4$. Recovering $q$
separately requires the $C^*$-level, where by Proposition 1.1 of Paper XV the point spectrum
$\Sp_p(\sigma)=q^{d_S\Z}$ is available and pins $q^{d_S}$ directly.
\end{remark}

\section{The Weil--Chebotarev obstruction}

\begin{proposition}[Corrected bound]\label{prop:weilbound}
Let $\chi$ be a nontrivial character of $G_S$ that does \emph{not} factor through the
constant-field tower; write $\chi=z^{\deg}\cdot\chi_0$ with $z$ a root of unity and $\chi_0$
geometric with conductor $\ff_\chi$, and put $B_\chi=2g-2+\deg\ff_\chi$
\cite{Stichtenoth}. Then $L(T,\chi)=L(zT,\chi_0)=\prod_{j=1}^{B_\chi}(1-\alpha_jT)$ with
$|\alpha_j|=q^{1/2}$ by Weil \cite{Weil}, and
\[
\Bigl|\sum_{\deg P=n}\chi(\Frob_P)\Bigr|\ \le\ \frac{B_\chi}{n}\,q^{n/2}+\frac{4g+4}{n}q^{n/2}
\qquad(n\ge1).
\]
\end{proposition}

\begin{proof}
Taking logarithmic derivatives of $L(T,\chi)$ gives
$\sum_{\deg P\mid n}\deg P\cdot\chi(\Frob_P)^{n/\deg P}=-\sum_j\alpha_j^n$, of modulus at
most $B_\chi q^{n/2}$. The terms with $\deg P=m<n$ are bounded in absolute value by
$m\,a_m\le\#C(\F_{q^m})\le q^m+1+2gq^{m/2}\le(2g+2)q^m$, and
$\sum_{m\le n/2}(2g+2)q^m\le(2g+2)\tfrac{q}{q-1}q^{n/2}\le(4g+4)q^{n/2}$; removing them and
dividing by $n$ gives the statement.
\end{proof}

\begin{remark}[The proposed bound]\label{rem:proposedbound}
The bound $2g_Lq^{n/2}/n+(g_L+1)q^{n/2}$ is true but its second term has no factor $1/n$ and
therefore dominates the first for every $n\ge1$, making the whole estimate weaker than the
trivial $O(q^{n/2})$ that the main term already provides. The correct order is
$O_\chi(q^{n/2}/n)$, as in Proposition~\ref{prop:weilbound}.
\end{remark}

\begin{remark}[Constant field extensions need separate treatment]\label{rem:constfield}
Weil's bound does not apply to the characters $\chi$ factoring through
$\Gal(\F_{q^m}/\F_q)$, for which $\chi(\Frob_P)=\zeta^{\deg P}$ with $\zeta$ a primitive
$m$-th root of unity: the corresponding $L$-function is $Z_C(\zeta T)$, which has poles
rather than the required polynomial form. These are the function field counterpart of the
class-group characters of \cite[\S4.1]{Main}, and they must be handled directly. That is
elementary: $D_\chi=\{P:\deg P\not\equiv0\ (\mathrm{mod}\ m)\}$, and
$\sum_{P\in D_\chi}q^{-\beta\deg P}=\sum_{n\not\equiv0\ (m)}a_nq^{-\beta n}$ diverges for
$\beta\le1$, since $a_n\sim q^n/n$ on the nonzero residue classes as well
(Corollary~\ref{cor:betac}).
\end{remark}

\begin{theorem}[Unconditional uniqueness]\label{thm:uniqueness}
Let $S=\cP_K$. For every nontrivial $\chi\in\widehat{G_S}$,
\[
\sum_{P}q^{-\beta\deg P}\bigl|1-\chi(\sigma_P)\bigr|^2=\infty\qquad(0<\beta\le1),
\]
so $\Xi_\beta=\{1\}$ and the $\mathrm{KMS}_\beta$ state is unique for $0<\beta\le1$,
unconditionally and with explicit constants. In particular the unramified class-group
obstruction of \cite[\S4.1]{Main} has no analogue at $S=\cP_K$: the characters pulled back
from $\operatorname{Pic}(C)$ have detecting sets of positive density.
\end{theorem}

\begin{proof}
Since $|1-w|^2=2-2\operatorname{Re}w$ on the unit circle,
\[
\sum_{\deg P=n}\bigl|1-\chi(\Frob_P)\bigr|^2
=2a_n-2\operatorname{Re}\sum_{\deg P=n}\chi(\Frob_P).
\]
If $\chi$ does not factor through the
constant-field tower, Proposition~\ref{prop:weilbound} bounds the second term by
$2(B_\chi+4g+4)q^{n/2}/n$, so the sum is $\ge a_n$ for $n$ large and
$\sum_nq^{-\beta n}a_n=\infty$ for $\beta\le1$ by Corollary~\ref{cor:betac}. If
$\chi=\zeta^{\deg}$ is a constant-field character, Remark~\ref{rem:constfield} applies.
These two cases are exhaustive, $\widehat{G_S}$ being torsion. Hence $\chi\notin\Xi_\beta$
for every nontrivial $\chi$ and $0<\beta\le1$, and uniqueness follows from the criterion of
\cite[Thm.~3.9]{Main}, whose Kakutani--martingale proof uses only the product structure of
the measures and applies verbatim over a function field. For the last sentence: a character
of $\operatorname{Pic}(C)$ pulled back to $G_S$ is $\zeta^{\deg}$ times a character of the
finite group $\operatorname{Pic}^0(C)$; the latter is geometric and everywhere unramified,
with $B_\chi=2g-2$, so it falls under the first case, and its detecting set has density
bounded below explicitly in $g$.
\end{proof}

\begin{remark}[What is gained over the number field case]\label{rem:gain}
The corresponding statement over $\Q$ \cite[Thm.~3.9]{Main} relies on the prime ideal theorem
and Chebotarev density with their ineffective error terms, Siegel--Walfisz and the possibility
of Landau--Siegel zeros. Here Weil's theorem is available and every constant is explicit in
$g$ and $\deg\ff_\chi$. Nothing is conjectural.
\end{remark}

\section{Unconditional sparse sets and closed phases}

\begin{theorem}[Closed low-temperature phase, no sieve]\label{thm:sparse}
Choose, for each $n$, any $A_n=\lfloor q^n/n^2\rfloor$ closed points of degree $n$; this is
possible for $n$ large since $a_n\sim q^n/n\gg q^n/n^2$. Let $S_{\mathrm{geom}}$ be their
union. Then
\[
\sum_{P\in S_{\mathrm{geom}}}q^{-\beta\deg P}\asymp\sum_n\frac{q^{n(1-\beta)}}{n^2},
\]
so $\beta_c(S_{\mathrm{geom}})=1$ and the series \emph{converges at} $\beta=1$, with value
$\asymp\zeta(2)=\pi^2/6$. Consequently the symmetry-broken phase is the closed half-line
$[1,\infty)$, exactly as in \cite[Thm.~5.5]{Main}, and the construction uses no sieve
hypothesis of any kind.
\end{theorem}

\begin{proof}
The displayed asymptotic is immediate; $\sum_nn^{-2}<\infty$ gives convergence at $\beta=1$
and $\sum_nq^{n(1-\beta)}/n^2=\infty$ for $\beta<1$ gives $\beta_c=1$. Closedness of the
phase is \cite[Thm.~5.5]{Main}, whose hypothesis is exactly convergence at the critical
exponent.
\end{proof}

\begin{remark}[Numerical check]\label{rem:numsparse}
With $q=2$, $\sum_{n\le200}2^{n(1-\beta)}/n^2$ equals $2.964$, $1.640$, $1.493$ at
$\beta=0.95,1.00,1.05$; the middle value is the truncation of $\pi^2/6=1.6449$.
\end{remark}

\begin{remark}[Comparison]\label{rem:compare}
Over $\Q$ the closed phase of \cite[Thm.~5.5]{Main} was obtained for the Sophie Germain
primes and required Brun's sieve; the corresponding statement for the twin primes remains
tied to the Hardy--Littlewood conjecture. Over a function field the point counts are known
exactly, so a set of any prescribed density can simply be selected. The gain is not that a
hard theorem becomes easy but that the question stops being about existence: one chooses the
set rather than proving that one exists.
\end{remark}

\section{Dictionary and assessment}

\[
\begin{array}{lll}
\text{Number field }K & \text{Function field }\F_q(C) & \text{Reference}\\\hline
\log\Ns\pp\ \text{never commensurable} & \text{always commensurable} & \text{\cite[Lem.~4.1]{SeriesI}, Lem.~3.1}\\
\sigma\ \text{never periodic} & \text{periodic},\ T_0=2\pi/(d_S\log q) & \text{Prop.~1.1}\\
\beta_c\ \text{depends on }S & \beta_c=1\ \text{for }S=\cP_K & \text{Cor.~2.2}\\
\text{type }\mathrm{III}_1\ \text{generic} & \textbf{never }\mathrm{III}_1 & \text{Thm.~3.3}\\
M\cong R_\infty,\ \beta\ \text{forgotten} & M\cong R_\lambda\ \text{recovers}\ \beta d_S\log q & \text{Thm.~4.1}\\
\text{Chebotarev ineffective} & \text{Weil bound, explicit} & \text{Prop.~5.1}\\
\text{sparse sets need sieves} & \text{sets are chosen, not found} & \text{Thm.~6.1}\\
\text{unramified obstruction (\cite[\S4.1]{Main})} & \text{none at }S=\cP_K\text{: positive density} & \text{Thm.~5.4}
\end{array}
\]

The transfer succeeds and, unlike the transfer to arithmetic topology of Paper VIII,
it succeeds with content: every analytic difficulty of the number field theory becomes a
theorem, and one structural feature genuinely inverts. That inversion, total commensurability
of the log-norms, is the whole story. It makes the flow periodic, it makes $\Gamma_\beta$
discrete, it forbids type $\mathrm{III}_1$, and it thereby causes the von Neumann algebra to
retain the temperature that Paper IX showed is destroyed over number fields.

We should be exact about the size of that last claim. The factor recovers one real number,
$\beta d_S\log q$, and by Remark~\ref{rem:ambiguity} not even $q$ and $\beta$ separately. It
does not recover the curve, its genus, or the obstruction group; the collapse of
Paper IX is weakened, not reversed. But it is the first weakening, and it comes from
the arithmetic of the ground field rather than from any additional structure imposed on the
system.

Three questions. Which of the three cases of Theorem~\ref{thm:neverIII1} actually occurs, and
for which $S$? The natural guess is $\mathrm{III}_{q^{-\beta d_S}}$ throughout, but
determining $\Sigma_\beta$ exactly, as Paper I does over number fields
(\cite[Prop.~2.2]{SeriesI} for the bound via the extended group, \cite[Thm.~4.2]{SeriesI}
for an exact determination), requires computing the extended obstruction group here and we
have not done it. Second, does the periodicity of $\sigma$ --- so that $\cA_{C,S}$ carries a
circle action rather than a flow --- bring the classification of Paper XIII within reach of
the equivariant Kirchberg theory for compact groups, which is far better developed than for
$\R$ (its Remark 5.1)? Third, the Weil bound gives effective control on
$\Xi_\beta$ that has no number field counterpart; whether it yields a realization theorem for
prescribed chains $\{\Xi_\beta\}$, the function field analogue of
\cite[Thm.~5.1]{SeriesI}, is open and looks tractable.

\begin{thebibliography}{9}
\bibitem{Rosen} M. Rosen, \emph{Number Theory in Function Fields}, Graduate Texts in Math. 210, Springer, 2002.
\bibitem{Stichtenoth} H. Stichtenoth, \emph{Algebraic Function Fields and Codes}, 2nd ed., Graduate Texts in Math. 254, Springer, 2009.
\bibitem{Weil} A. Weil, \emph{Sur les courbes alg\'ebriques et les vari\'et\'es qui s'en d\'eduisent}, Hermann, Paris, 1948.
\bibitem{Main} R. Chen, \emph{KMS state uniqueness and phase transitions in localized Bost--Connes systems}, preprint, 2026, \newline\url{https://doi.org/10.5281/zenodo.22152101}.
\bibitem{SeriesI} R. Chen, \emph{Localized Bost--Connes systems I: factor type and the spectrum of transitions}, preprint, 2026, DOI \href{https://doi.org/10.5281/zenodo.22160827}{10.5281/zenodo.22160827}.
\end{thebibliography}

\end{document}
